\section{Details of Experiments on QM9}
\label{appendix:sec:qm9}



\begin{table}[]
\centering
\resizebox{\ifdim\width>\textwidth\textwidth\else\width\fi}{!}{
\begin{tabular}{ll}
%\cline{7-11}
\toprule[1.2pt]
Hyper-parameters & Value or description
 \\
\midrule[1.2pt]
Optimizer & AdamW \\
Learning rate scheduling & Cosine learning rate with linear warmup \\
Warmup epochs & $5$ \\
%Maximum learning rate & $1.5 \times 10 ^{-4}$ for $G$, $H$, $U$ and $U_0$, and $5 \times 10 ^{-4}$ for others \\
Maximum learning rate & $1.5 \times 10 ^{-4}$, $5 \times 10 ^{-4}$ \\
%Batch size & $64$ for $G$, $H$, $U$ and $U_0$, and $128$ for others \\
Batch size & $64$, $128$ \\
%Number of epochs & $600$ for $G$, $H$, $U$ and $U_0$, and $300$ for others \\
Number of epochs & $300$, $600$ \\
%Weight decay & $0$ for $G$, $H$, $U$ and $U_0$, and $5 \times 10^{-3}$ for others \\
Weight decay & $0$, $5 \times 10^{-3}$ \\
%Dropout rate & $0.0$ for $G$, $H$, $U$ and $U_0$, $0.1$ for $R^{2}$, and $0.2$ for others \\
Dropout rate & $0.0$, $0.1$, $0.2$ \\
% ===============================
%   Add this back in the futrure
% ===============================
%Dropout rate & $0.1$, $0.2$ \\
\midrule[0.6pt]
Cutoff radius ($\angstrom$) & $5$ \\
Number of radial bases & $128$ for Gaussian radial basis, $8$ for radial bessel basis \\
Hidden sizes of radial functions & $64$ \\
Number of hidden layers in radial functions & $2$ \\

\midrule[0.6pt]

\multicolumn{2}{c}{Equiformer} \\
\\
Number of Transformer blocks & $6$ \\
Embedding dimension $d_{embed}$ & $[(128, 0), (64, 1), (32, 2)]$ \\
Spherical harmonics embedding dimension $d_{sh}$ & $[(1, 0), (1, 1), (1, 2)]$ \\
Number of attention heads $h$ & $4$ \\
Attention head dimension $d_{head}$ & $[(32, 0), (16, 1), (8, 2)]$\\
Hidden dimension in feed forward networks $d_{ffn}$ & $[(384, 0), (192, 1) , (96, 2)]$ \\
Output feature dimension $d_{feature}$ & $[(512, 0)]$\\

\midrule[0.6pt]

\multicolumn{2}{c}{$E(3)$-Equiformer} \\
\\
Number of Transformer blocks & $6$ \\
Embedding dimension $d_{embed}$ & $[(128, 0, e), (32, 0, o), (32, 1, e), (32, 1, o), (16, 2, e), (16, 2, o)]$ \\
Spherical harmonics embedding dimension $d_{sh}$ & $[(1, 0, e), (1, 1, o), (1, 2, e)]$ \\
Number of attention heads $h$ & $4$ \\
Attention head dimension $d_{head}$ & $[(32, 0, e), (8, 0, o), (8, 1, e), (8, 1, o), (4, 2, e), (4, 2, o)]$ \\
Hidden dimension in feed forward networks $d_{ffn}$ & $[(384, 0, e), (96, 0, o), (96, 1, e), (96, 1, o), (48, 2, e), (48, 2, o)]$ \\
Output feature dimension $d_{feature}$ & $[(512, 0, e)]$\\
\bottomrule[1.2pt]
\end{tabular}
}
\vspace{-2mm}
\caption{\textbf{Hyper-parameters for QM9 dataset.}
We denote $C_L$ type-$L$ vectors as $(C_L, L)$ and $C_{(L, p)}$ type-$(L, p)$ vectors as $(C_{(L, p)}, L, p)$ and use brackets to represent concatenations of vectors.
}
\label{appendix:tab:qm9}
\end{table}


\subsection{Training Details}
\label{appendix:subsec:qm9_training_details}
%%%We normalize ground truth by subtracting mean and dividing by standard deviation.
%%%For the task of $U$, $U_0$, $G$, and $H$, where single-atom reference values are available, we subtract those reference values from ground truth before normalizing. 
We use the same data partition as TorchMD-NET~\citep{torchmd_net}.
For the task of $U$, $U_0$, $G$, and $H$, where single-atom reference values are available, we subtract those reference values from ground truth. 
For other tasks, we normalize ground truth by subtracting mean and dividing by standard deviation.

We train Equiformer with 6 blocks with $L_{max} = 2$.
%We choose Gaussian radial basis~\cite{schnet, spinconv, gemnet, graphormer_3d} for the first six tasks in Table~\ref{tab:qm9_result} and radial Bessel basis for the others.
We choose Gaussian radial basis~\citep{schnet, spinconv, gemnet, graphormer_3d} for the first six tasks in Table~\ref{tab:qm9_result} and radial Bessel basis~\citep{dimenet, dimenet_pp} for the others.
% ===============================
%   Add these back in the future
% ===============================
We apply dropout~\citep{dropout} to attention weights $a_{ij}$.
The dropout rate is $0.0$ for the tasks of $G$, $H$, $U$ and $U_0$, is $0.1$ for the task of $R^2$ and is $0.2$ for others.
Since the tasks of $G$, $H$, $U$, and $U_0$ require longer training, we use slightly different hyper-parameters.
The number of epochs is $600$ for the tasks of $G$, $H$, $U$, and $U_0$ and is $300$ for others.
The learning rate is $1.5 \times 10 ^{-4}$ for the tasks of $G$, $H$, $U$, and $U_0$ and is $5 \times 10^{-4}$ for others.
The batch size is $64$ for the tasks of $G$, $H$, $U$, and $U_0$ and is $128$ for others.
The weight decay is $0$ for the tasks of $G$, $H$, $U$, and $U_0$ and is $5 \times 10^{-3}$ for others.
Table~\ref{appendix:tab:qm9} summarizes the hyper-parameters for the QM9 dataset.
The detailed description of architectural hyper-parameters can be found in Sec.~\ref{appendix:subsec:equiformer_architecture}. 

We use one A6000 GPU with 48GB to train each model and summarize the computational cost of training for one epoch as follows.
Training $E(3)$-Equiformer in Table~\ref{appendix:tab:e3_qm9} for one epoch takes about $16.3$ minutes.
The time of training Equiformer, Equiformer with linear messages (indicated by Index $2$ in Table~\ref{tab:qm9_ablation}), and Equiformer with linear messages and dot product attention (indicated by Index $3$ in Table~\ref{tab:qm9_ablation}) for one epoch is $12.1$ minutes, $7.2$ minutes and $7.8$ minutes, respectively.


%We set cutoff radius to be $5 \angstrom$ and represent relative distances with Gaussian radial basis~\cite{schnet, spinconv, gemnet, graphormer_3d} with dimension $128$.
%Please refer to appendix for further details on Equiformer architecture.
%\todo{We apply dropout~\cite{dropout} with probability $0.2$ to attention weights}. 
%We use AdamW optimizer~\cite{adamw} with weight decay $5 \times 10^{-3}$ and use cosine learning rate with linear warmup, where the maximum learning rate is $5 \times 10 ^{-4}$ and warmup epoch is $5$.
%The batch size is $128$ and the number of epochs is $300$.


\subsection{Comparison between $SE(3)$ and $E(3)$ Equivariance}
\label{appendix:subsec:e3_qm9}
We train two versions of Equiformers, one with $SE(3)$-equivariant features denoted as ``Equiformer'' and the other with $E(3)$-equivariant features denoted as ``$E(3)$-Equiformer'', and we compare them in Table~\ref{appendix:tab:e3_qm9}.
%Including equivariance to inversion further improves the performance on QM9 dataset.
As for Table~\ref{tab:qm9_result}, we compare ``Equiformer'' with other works since most of them do not include equivariance to inversion.

\begin{table}[th]
\centering
\resizebox{\ifdim\width>\textwidth\textwidth\else\width\fi}{!}{
\begin{tabular}{llccccccccc}
%\cline{7-11}
\toprule[1.2pt]
& Task & $\alpha$ & $\Delta \varepsilon$ & $\varepsilon_{\text{HOMO}}$ & $\varepsilon_{\text{LUMO}}$ & $\mu$ & $C_{\nu}$ & Training time & Number of \\
%Methods & Units & bohr$^3$ & meV & meV & meV & D & cal/mol K & (minutes/epoch) & parameters \\
Methods & Units & $a_0^3$ & meV & meV & meV & D & cal/mol K & (minutes/epoch) & parameters \\
\midrule[1.2pt]


Equiformer & & .046 & 30 & 15 & 14 & .011 & .023 & 12.1 & 3.53M \\
$E(3)$-Equiformer & & .045 & 30 & 15 & 14 & .012 & .023 & 16.3 & 3.28M\\
%\midrule[0.6pt]
%Equiformer + NAT & 0.4737 & 0.7245 & 0.4862 & 0.6493 & 0.5834 & 6.01 & 2.45 & 5.41 & 2.71 & \\
\bottomrule[1.2pt]

\end{tabular}
}
\vspace{-2mm}
\caption{\textbf{Ablation study of $SE(3)$/$E(3)$ equivariance on QM9 testing set.}  
``Equiformer'' operates on $SE(3)$-equivariant features while ``$E(3)$-Equiformer'' uses $E(3)$-equivariant features.
Including inversion achieves similar performance.
}
\vspace{-2.5mm}
\label{appendix:tab:e3_qm9}
\end{table}


\subsection{Comparison of Training Time and Numbers of Parameters}
\label{appendix:subsec:qm9_training_time_comparison}
We compare training time and numbers of parameters between SEGNN~\citep{segnn}, TorchMD-NET~\citep{torchmd_net} and Equiformer and summarize the results in Table~\ref{appendix:tab:qm9_training_time}.
Training Equiformer for $300$ epochs and for $600$ epochs takes $61$ and $122$ GPU-hours, respectively.

Compared to SEGNN, which is written with the same \texttt{e3nn} library~\citep{e3nn}, Equiformer with MLP attention and non-linear message is faster.
Although Equiformer has more channels and more parameters, the training time is comparable. 
The reasons are as follows.
Equiformer uses more efficient depth-wise tensor products (DTP), where one output channel depends on only one input channel. 
SEGNN uses more compute-intensive fully connected tensor products (FCTP), where one output channel depends on all input channels. 
Besides, SEGNN uses 4 FCTPs in each message passing block while Equiformer uses only 2 DTPs in each block. 

Compared to TorchMD-NET, which is trained for $3000$ epochs, Equiformer achieves competitve results after trained for $300$ or $600$ epochs.
Equiformer takes more time for each epoch since Equiformer uses more expressive non-linear messages, which compared to linear messages used in other equivariant Transformers, doubles the number of tensor products and therefore almost doubles the training time.
Moreover, Equiformer incorporates tensors of higher degrees (e.g., $L_{max}=2$), which improves performance but slows down the training.



\begin{table}[t]
\begin{adjustwidth}{0cm}{0cm}
\centering
\resizebox{\ifdim\width>\textwidth\textwidth\else\width\fi}{!}{
\begin{tabular}{lccc}
%\cline{7-11}
\toprule[1.2pt]
Methods & Number of parameters & Training time (GPU-hours) \\
\midrule[1.2pt]

SEGNN~\citep{segnn} & 1.03M & 81 \\
TorchMD-NET~\citep{torchmd_net} & 6.86M & 92 \\
\midrule[0.6pt]
Equiformer & 3.53M & 61 \\

\bottomrule[1.2pt]

\end{tabular}
}
\vspace{-3mm}
\end{adjustwidth}
\caption{\textbf{Comparison of training time and numbers of parameters for QM9 dataset.}  
%\note{Dropout rate = 0.1.}
%%%$\dagger$ denotes using different training, validation, testing data partitions. % as mentioned in SEGNN~\cite{segnn}.
%$\dagger$ denotes using different data partitions.
%$\ddagger$ denotes results from SE(3)-Transformer~\cite{se3_transformer}.
}
\label{appendix:tab:qm9_training_time}
\end{table}